% Seminar 4 - Pagina 5. Fragment pentru
% % Seminar 4 - Pagina 5. Fragment pentru
% % Seminar 4 - Pagina 5. Fragment pentru
% % Seminar 4 - Pagina 5. Fragment pentru
% \input{seminars/pages/seminar4-pagina5-ro}; fără preambul sau pagină nouă.
% NOTE EDITORIALE:
% - Prima observație reprezintă un contraexemplu la o analogie de tip
%   Cantor-Bernstein pentru grupuri: fiecare grup poate fi izomorf cu
%   un subgrup al celuilalt, fără ca grupurile să fie izomorfe.
%   Săgeata urmată de G_1 neizomorf cu G_2 este păstrată în acest sens,
%   nu ca implicație general valabilă.
% - Manuscrisul scrie (q_0,q_1,...,q_n), dar aplicația de deplasare
%   phi((q_0,...,q_n))=(0,q_0,...) este izomorfismul dorit doar dacă
%   șirul de coordonate continuă la infinit. Ambiguitatea de indexare
%   NU este corectată tacit în formulele afișate. La n finit fix,
%   Q^(n+1) nu este izomorf cu 0 x Q^n.
% - Manuscrisul nu demonstrează neizomorfismul G_1, G_2. Pentru șiruri
%   infinite, G_1 este divizibil, iar G_2, cu prima coordonată întreagă,
%   nu este divizibil.
% - Ultima notă indică ordinele 4,5,6,7, dar enumeră explicit doar
%   grupurile de ordine 4,5,6; cazul 7 nu este completat în manuscris.
% - Înaintea liniei Z_6/S_3 există o corectură cu cerneală; au fost
%   transcrise numai denumirile lizibile. Urmele de pe verso sunt omise.

\subsection*{Seminar 4 - Pagina 5}

\noindent Fie $G_1,G_2$ grupuri. Manuscrisul consideră situația
\[
\begin{gathered}
 G_1\text{ izomorf cu un subgrup al lui }G_2,\\
 G_2\text{ izomorf cu un subgrup al lui }G_1,
\end{gathered}
\qquad\text{dar}\qquad G_1\not\cong G_2.
\]

\medskip
\noindent Definim
\[
 G_1:=\bigl\{x=(q_0,q_1,\ldots,q_n)\mid q_i\in\mathbb Q\bigr\}.
\]
Atunci $(G_1,+)$ este grup. Pentru
$y=(u_0,u_1,\ldots,u_n)$,
\[
 x+y=(q_0+u_0,\ldots,q_n+u_n).
\]

\noindent De asemenea,
\[
 \begin{aligned}
 G_2:=\bigl\{\,&x=(q_0,q_1,\ldots,q_n)\mid
     q_0\in\mathbb Z,\quad q_i\in\mathbb Q\ 
     (1\leq i\leq n)\,\bigr\}\subseteq G_1.
 \end{aligned}
\]
$(G_2,+)$ este grup și $G_2\leq G_1$
($G_2$ subgrup al lui $G_1$).

\medskip
\noindent Fie
\[
 H:=\bigl\{x=(0,q_1,\ldots,q_n)\mid
     q_i\in\mathbb Q\ (1\leq i\leq n)\bigr\}.
\]
Atunci $H\leq G_2$. În manuscris este dată aplicația
\[
 \begin{aligned}
 \varphi:G_1&\longrightarrow H,\\
 \varphi\bigl((q_0,\ldots,q_n)\bigr)&=(0,q_0,\ldots),
 \end{aligned}
\]
și se notează că $\varphi$ este izomorfism de grupuri.

\medskip
\noindent\textbf{Concluzie (după manuscris).}
$G_1$ și $G_2$ nu sunt grupuri izomorfe.

\medskip
\noindent\underline{Grupuri de ordin mic.}
Manuscrisul enumeră $|G|\in\{4,5,6,7\}$ și consemnează
\[
\begin{array}{r@{\quad\Longrightarrow\quad}l}
 |G|=4 & G\cong\mathbb Z_4\quad\text{sau}\quad
               \mathbb Z_2\times\mathbb Z_2,\\[2pt]
 |G|=5 & G\cong\mathbb Z_5,\\[2pt]
 |G|=6 & G\cong\mathbb Z_6\quad\text{sau}\quad S_3.
\end{array}
\]
; fără preambul sau pagină nouă.
% NOTE EDITORIALE:
% - Prima observație reprezintă un contraexemplu la o analogie de tip
%   Cantor-Bernstein pentru grupuri: fiecare grup poate fi izomorf cu
%   un subgrup al celuilalt, fără ca grupurile să fie izomorfe.
%   Săgeata urmată de G_1 neizomorf cu G_2 este păstrată în acest sens,
%   nu ca implicație general valabilă.
% - Manuscrisul scrie (q_0,q_1,...,q_n), dar aplicația de deplasare
%   phi((q_0,...,q_n))=(0,q_0,...) este izomorfismul dorit doar dacă
%   șirul de coordonate continuă la infinit. Ambiguitatea de indexare
%   NU este corectată tacit în formulele afișate. La n finit fix,
%   Q^(n+1) nu este izomorf cu 0 x Q^n.
% - Manuscrisul nu demonstrează neizomorfismul G_1, G_2. Pentru șiruri
%   infinite, G_1 este divizibil, iar G_2, cu prima coordonată întreagă,
%   nu este divizibil.
% - Ultima notă indică ordinele 4,5,6,7, dar enumeră explicit doar
%   grupurile de ordine 4,5,6; cazul 7 nu este completat în manuscris.
% - Înaintea liniei Z_6/S_3 există o corectură cu cerneală; au fost
%   transcrise numai denumirile lizibile. Urmele de pe verso sunt omise.

\subsection*{Seminar 4 - Pagina 5}

\noindent Fie $G_1,G_2$ grupuri. Manuscrisul consideră situația
\[
\begin{gathered}
 G_1\text{ izomorf cu un subgrup al lui }G_2,\\
 G_2\text{ izomorf cu un subgrup al lui }G_1,
\end{gathered}
\qquad\text{dar}\qquad G_1\not\cong G_2.
\]

\medskip
\noindent Definim
\[
 G_1:=\bigl\{x=(q_0,q_1,\ldots,q_n)\mid q_i\in\mathbb Q\bigr\}.
\]
Atunci $(G_1,+)$ este grup. Pentru
$y=(u_0,u_1,\ldots,u_n)$,
\[
 x+y=(q_0+u_0,\ldots,q_n+u_n).
\]

\noindent De asemenea,
\[
 \begin{aligned}
 G_2:=\bigl\{\,&x=(q_0,q_1,\ldots,q_n)\mid
     q_0\in\mathbb Z,\quad q_i\in\mathbb Q\ 
     (1\leq i\leq n)\,\bigr\}\subseteq G_1.
 \end{aligned}
\]
$(G_2,+)$ este grup și $G_2\leq G_1$
($G_2$ subgrup al lui $G_1$).

\medskip
\noindent Fie
\[
 H:=\bigl\{x=(0,q_1,\ldots,q_n)\mid
     q_i\in\mathbb Q\ (1\leq i\leq n)\bigr\}.
\]
Atunci $H\leq G_2$. În manuscris este dată aplicația
\[
 \begin{aligned}
 \varphi:G_1&\longrightarrow H,\\
 \varphi\bigl((q_0,\ldots,q_n)\bigr)&=(0,q_0,\ldots),
 \end{aligned}
\]
și se notează că $\varphi$ este izomorfism de grupuri.

\medskip
\noindent\textbf{Concluzie (după manuscris).}
$G_1$ și $G_2$ nu sunt grupuri izomorfe.

\medskip
\noindent\underline{Grupuri de ordin mic.}
Manuscrisul enumeră $|G|\in\{4,5,6,7\}$ și consemnează
\[
\begin{array}{r@{\quad\Longrightarrow\quad}l}
 |G|=4 & G\cong\mathbb Z_4\quad\text{sau}\quad
               \mathbb Z_2\times\mathbb Z_2,\\[2pt]
 |G|=5 & G\cong\mathbb Z_5,\\[2pt]
 |G|=6 & G\cong\mathbb Z_6\quad\text{sau}\quad S_3.
\end{array}
\]
; fără preambul sau pagină nouă.
% NOTE EDITORIALE:
% - Prima observație reprezintă un contraexemplu la o analogie de tip
%   Cantor-Bernstein pentru grupuri: fiecare grup poate fi izomorf cu
%   un subgrup al celuilalt, fără ca grupurile să fie izomorfe.
%   Săgeata urmată de G_1 neizomorf cu G_2 este păstrată în acest sens,
%   nu ca implicație general valabilă.
% - Manuscrisul scrie (q_0,q_1,...,q_n), dar aplicația de deplasare
%   phi((q_0,...,q_n))=(0,q_0,...) este izomorfismul dorit doar dacă
%   șirul de coordonate continuă la infinit. Ambiguitatea de indexare
%   NU este corectată tacit în formulele afișate. La n finit fix,
%   Q^(n+1) nu este izomorf cu 0 x Q^n.
% - Manuscrisul nu demonstrează neizomorfismul G_1, G_2. Pentru șiruri
%   infinite, G_1 este divizibil, iar G_2, cu prima coordonată întreagă,
%   nu este divizibil.
% - Ultima notă indică ordinele 4,5,6,7, dar enumeră explicit doar
%   grupurile de ordine 4,5,6; cazul 7 nu este completat în manuscris.
% - Înaintea liniei Z_6/S_3 există o corectură cu cerneală; au fost
%   transcrise numai denumirile lizibile. Urmele de pe verso sunt omise.

\subsection*{Seminar 4 - Pagina 5}

\noindent Fie $G_1,G_2$ grupuri. Manuscrisul consideră situația
\[
\begin{gathered}
 G_1\text{ izomorf cu un subgrup al lui }G_2,\\
 G_2\text{ izomorf cu un subgrup al lui }G_1,
\end{gathered}
\qquad\text{dar}\qquad G_1\not\cong G_2.
\]

\medskip
\noindent Definim
\[
 G_1:=\bigl\{x=(q_0,q_1,\ldots,q_n)\mid q_i\in\mathbb Q\bigr\}.
\]
Atunci $(G_1,+)$ este grup. Pentru
$y=(u_0,u_1,\ldots,u_n)$,
\[
 x+y=(q_0+u_0,\ldots,q_n+u_n).
\]

\noindent De asemenea,
\[
 \begin{aligned}
 G_2:=\bigl\{\,&x=(q_0,q_1,\ldots,q_n)\mid
     q_0\in\mathbb Z,\quad q_i\in\mathbb Q\ 
     (1\leq i\leq n)\,\bigr\}\subseteq G_1.
 \end{aligned}
\]
$(G_2,+)$ este grup și $G_2\leq G_1$
($G_2$ subgrup al lui $G_1$).

\medskip
\noindent Fie
\[
 H:=\bigl\{x=(0,q_1,\ldots,q_n)\mid
     q_i\in\mathbb Q\ (1\leq i\leq n)\bigr\}.
\]
Atunci $H\leq G_2$. În manuscris este dată aplicația
\[
 \begin{aligned}
 \varphi:G_1&\longrightarrow H,\\
 \varphi\bigl((q_0,\ldots,q_n)\bigr)&=(0,q_0,\ldots),
 \end{aligned}
\]
și se notează că $\varphi$ este izomorfism de grupuri.

\medskip
\noindent\textbf{Concluzie (după manuscris).}
$G_1$ și $G_2$ nu sunt grupuri izomorfe.

\medskip
\noindent\underline{Grupuri de ordin mic.}
Manuscrisul enumeră $|G|\in\{4,5,6,7\}$ și consemnează
\[
\begin{array}{r@{\quad\Longrightarrow\quad}l}
 |G|=4 & G\cong\mathbb Z_4\quad\text{sau}\quad
               \mathbb Z_2\times\mathbb Z_2,\\[2pt]
 |G|=5 & G\cong\mathbb Z_5,\\[2pt]
 |G|=6 & G\cong\mathbb Z_6\quad\text{sau}\quad S_3.
\end{array}
\]
; fără preambul sau pagină nouă.
% NOTE EDITORIALE:
% - Prima observație reprezintă un contraexemplu la o analogie de tip
%   Cantor-Bernstein pentru grupuri: fiecare grup poate fi izomorf cu
%   un subgrup al celuilalt, fără ca grupurile să fie izomorfe.
%   Săgeata urmată de G_1 neizomorf cu G_2 este păstrată în acest sens,
%   nu ca implicație general valabilă.
% - Manuscrisul scrie (q_0,q_1,...,q_n), dar aplicația de deplasare
%   phi((q_0,...,q_n))=(0,q_0,...) este izomorfismul dorit doar dacă
%   șirul de coordonate continuă la infinit. Ambiguitatea de indexare
%   NU este corectată tacit în formulele afișate. La n finit fix,
%   Q^(n+1) nu este izomorf cu 0 x Q^n.
% - Manuscrisul nu demonstrează neizomorfismul G_1, G_2. Pentru șiruri
%   infinite, G_1 este divizibil, iar G_2, cu prima coordonată întreagă,
%   nu este divizibil.
% - Ultima notă indică ordinele 4,5,6,7, dar enumeră explicit doar
%   grupurile de ordine 4,5,6; cazul 7 nu este completat în manuscris.
% - Înaintea liniei Z_6/S_3 există o corectură cu cerneală; au fost
%   transcrise numai denumirile lizibile. Urmele de pe verso sunt omise.

\subsection*{Seminar 4 - Pagina 5}

\noindent Fie $G_1,G_2$ grupuri. Manuscrisul consideră situația
\[
\begin{gathered}
 G_1\text{ izomorf cu un subgrup al lui }G_2,\\
 G_2\text{ izomorf cu un subgrup al lui }G_1,
\end{gathered}
\qquad\text{dar}\qquad G_1\not\cong G_2.
\]

\medskip
\noindent Definim
\[
 G_1:=\bigl\{x=(q_0,q_1,\ldots,q_n)\mid q_i\in\mathbb Q\bigr\}.
\]
Atunci $(G_1,+)$ este grup. Pentru
$y=(u_0,u_1,\ldots,u_n)$,
\[
 x+y=(q_0+u_0,\ldots,q_n+u_n).
\]

\noindent De asemenea,
\[
 \begin{aligned}
 G_2:=\bigl\{\,&x=(q_0,q_1,\ldots,q_n)\mid
     q_0\in\mathbb Z,\quad q_i\in\mathbb Q\ 
     (1\leq i\leq n)\,\bigr\}\subseteq G_1.
 \end{aligned}
\]
$(G_2,+)$ este grup și $G_2\leq G_1$
($G_2$ subgrup al lui $G_1$).

\medskip
\noindent Fie
\[
 H:=\bigl\{x=(0,q_1,\ldots,q_n)\mid
     q_i\in\mathbb Q\ (1\leq i\leq n)\bigr\}.
\]
Atunci $H\leq G_2$. În manuscris este dată aplicația
\[
 \begin{aligned}
 \varphi:G_1&\longrightarrow H,\\
 \varphi\bigl((q_0,\ldots,q_n)\bigr)&=(0,q_0,\ldots),
 \end{aligned}
\]
și se notează că $\varphi$ este izomorfism de grupuri.

\medskip
\noindent\textbf{Concluzie (după manuscris).}
$G_1$ și $G_2$ nu sunt grupuri izomorfe.

\medskip
\noindent\underline{Grupuri de ordin mic.}
Manuscrisul enumeră $|G|\in\{4,5,6,7\}$ și consemnează
\[
\begin{array}{r@{\quad\Longrightarrow\quad}l}
 |G|=4 & G\cong\mathbb Z_4\quad\text{sau}\quad
               \mathbb Z_2\times\mathbb Z_2,\\[2pt]
 |G|=5 & G\cong\mathbb Z_5,\\[2pt]
 |G|=6 & G\cong\mathbb Z_6\quad\text{sau}\quad S_3.
\end{array}
\]
